\documentclass{article}
\usepackage{amsmath,amssymb,amsthm}
\usepackage{algorithm}
\usepackage{algorithmic}
\usepackage{hyperref}
\usepackage{geometry}
\geometry{margin=1in}
\newtheorem{theorem}{Theorem}
\newtheorem{lemma}{Lemma}
\newtheorem{assumption}{Assumption}
\newcommand{\EE}{\mathbb{E}}
\newcommand{\RR}{\mathbb{R}}
\newcommand{\norm}[1]{\left\|#1\right\|}
\begin{document}

\section*{Cyclic Stochastic Gradient Descent (Cyclic SGD)}

\section*{Mathematical Formulation}
Let $f:\RR^{d}\rightarrow\RR$ be the (possibly non‑convex) objective function.  
At iteration $t\in\{0,1,\dots,T-1\}$ we draw a random sample $z_t$ and compute a stochastic gradient
\[
g_t \;:=\; \nabla f(w_t;z_t)\in\RR^{d}.
\]
The learning rate follows a \emph{triangular cyclic schedule} with period $2M$:
\[
\eta_t \;=\; \eta_{\min}+(\eta_{\max}-\eta_{\min})\,
\Bigl(1-\Bigl|\frac{(t\bmod 2M)}{M}-1\Bigr|\Bigr),\qquad
t=0,1,\dots
\tag{1}
\]
Thus $\eta_t$ linearly increases from $\eta_{\min}$ to $\eta_{\max}$ in the first $M$ steps,
and linearly decreases back to $\eta_{\min}$ in the next $M$ steps, repeating thereafter.

The update rule is
\[
w_{t+1}\;=\;w_t-\eta_t\,g_t .
\tag{2}
\]

\section*{Pseudocode}
\begin{algorithm}[H]
\caption{Cyclic SGD}
\begin{algorithmic}[1]
\REQUIRE Initial point $w_0\in\RR^{d}$, 
          minimum step $\eta_{\min}>0$, 
          maximum step $\eta_{\max}> \eta_{\min}$,
          cycle length $M\in\mathbb{N}$,
          total iterations $T$.
\FOR{$t=0$ \TO $T-1$}
    \STATE Compute learning rate $\eta_t$ by \eqref{1}.
    \STATE Sample data $z_t$ and evaluate stochastic gradient $g_t=\nabla f(w_t;z_t)$.
    \STATE Update $w_{t+1}=w_t-\eta_t g_t$.
\ENDFOR
\ENSURE Final iterate $w_T$ (or optionally the average $\bar w_T=\frac1T\sum_{t=0}^{T-1}w_t$).
\end{algorithmic}
\end{algorithm}

\section*{Dry Run Example}
Consider the one‑dimensional quadratic $f(w)=(w-3)^2$.
Hence $\nabla f(w)=2(w-3)$.  
Set $w_0=0$, $\eta_{\min}=0.05$, $\eta_{\max}=0.1$, $M=2$ (so the schedule is $0.05,0.1,0.1,0.05,\dots$).  
Assume a deterministic gradient (i.e. $g_t=\nabla f(w_t)$).

\begin{center}
\begin{tabular}{c|c|c|c}
$t$ & $\eta_t$ & $g_t=2(w_t-3)$ & $w_{t+1}=w_t-\eta_t g_t$ \\ \hline
0 & $0.05$ & $2(0-3)=-6$ & $0-0.05(-6)=0.30$ \\
1 & $0.10$ & $2(0.30-3)=-5.40$ & $0.30-0.10(-5.40)=0.84$ \\
2 & $0.10$ & $2(0.84-3)=-4.32$ & $0.84-0.10(-4.32)=1.272$ 
\end{tabular}
\end{center}
Corresponding objective values:
$f(w_0)=9$, $f(w_1)=7.29$, $f(w_2)=5.29$, $f(w_3)=3.96$.
The cyclic learning rates lead to a faster reduction than a constant $\eta=0.05$ would.

\newpage

\section*{Assumptions}
\begin{assumption}[Smoothness]\label{as:smooth}
$f$ is $L$‑smooth, i.e.
\[
f(u)\le f(v)+\langle\nabla f(v),u-v\rangle+\frac{L}{2}\norm{u-v}^2,
\qquad\forall\,u,v\in\RR^{d}.
\tag{A1}
\]
\end{assumption}
\begin{assumption}[Unbiased Stochastic Gradient]\label{as:unbiased}
For every $t$, conditional on $w_t$,
\[
\EE\!\left[g_t\mid w_t\right]=\nabla f(w_t).
\tag{A2}
\]
\end{assumption}
\begin{assumption}[Bounded Variance]\label{as:variance}
There exists $\sigma^2\ge0$ such that
\[
\EE\!\left[\norm{g_t-\nabla f(w_t)}^2\mid w_t\right]\le\sigma^2,
\qquad\forall t.
\tag{A3}
\]
\end{assumption}
\begin{assumption}[Learning‑rate Bounds]\label{as:lr}
The cyclic schedule satisfies
\[
0<\eta_{\min}\le\eta_t\le\eta_{\max}\quad\text{for all }t,
\qquad 
\eta_{\max}\le\frac{1}{L}.
\tag{A4}
\]
\end{assumption}

\section*{Main Convergence Theorem}
\begin{theorem}[Convergence of Cyclic SGD]\label{thm:main}
Under Assumptions \ref{as:smooth}–\ref{as:lr}, for any $T\ge1$ the iterates produced by
Cyclic SGD satisfy
\[
\frac{1}{T}\sum_{t=0}^{T-1}\EE\!\left[\norm{\nabla f(w_t)}^2\right]
\;\le\;
\underbrace{\frac{2\bigl(f(w_0)-f^{\star}\bigr)}{\eta_{\min} T}}_{\text{optimization term}}
\;+\;
\underbrace{L\,\eta_{\max}\,\sigma^2}_{\text{variance term}},
\tag{3}
\]
where $f^{\star}=\inf_{w}f(w)$.
Consequently,
\[
\min_{0\le t<T}\EE\!\left[\norm{\nabla f(w_t)}^2\right]
=O\!\Bigl(\tfrac{1}{\eta_{\min}T}+\eta_{\max}\Bigr).
\]
\end{theorem}

\section*{Theoretical Properties}
\begin{itemize}
\item \textbf{Stability.} Because $\eta_{\max}\le 1/L$, the update \eqref{2} never exceeds the
standard gradient‑descent stability bound for $L$‑smooth functions.
\item \textbf{Hyper‑parameter Sensitivity.}
The bound \eqref{3} shows a linear dependence on $\eta_{\max}$ (through the variance term)
and an inverse dependence on $\eta_{\min}$ (through the optimization term).  Hence,
choosing a large $\eta_{\max}$ accelerates early progress but amplifies noise; a larger
$\eta_{\min}$ reduces the $1/T$ term but may violate the stability condition.
\item \textbf{Comparison to Constant‑step SGD.}
If $\eta_{\min}=\eta_{\max}=\eta$, \eqref{3} reduces to the classic bound
$\frac{2(f(w_0)-f^{\star})}{\eta T}+L\eta\sigma^2$, confirming that the cyclic schedule
inherits the same $O(1/T)$ rate while allowing a larger maximal step without sacrificing
the $1/T$ term thanks to the presence of $\eta_{\min}$ in the denominator.
\end{itemize}

\section*{Discussion}
The theorem guarantees that, despite the non‑convex landscape,
the average squared gradient norm diminishes at the rate $O(1/T)$ up to a
steady‑state error proportional to $\eta_{\max}\sigma^2$.
By cycling the step size, practitioners can enjoy rapid descent early in each
cycle (large $\eta_{\max}$) while preserving a small long‑run error thanks to the
lower bound $\eta_{\min}$ that appears in the dominant $1/T$ term.

\newpage

\section*{Preliminaries (Definitions and Lemmas)}
\begin{lemma}[Smoothness inequality]\label{lem:smooth}
If $f$ is $L$‑smooth, then for any $u,v\in\RR^{d}$,
\[
f(u)\le f(v)+\langle\nabla f(v),u-v\rangle+\frac{L}{2}\norm{u-v}^2 .
\tag{L1}
\]
\end{lemma}
\begin{proof}
Standard consequence of $L$‑smoothness; see Assumption \ref{as:smooth}. \qedhere
\end{proof}

\begin{lemma}[Descent Lemma for a single step]\label{lem:descent}
For the update $w_{t+1}=w_t-\eta_t g_t$, we have
\[
\EE\!\left[f(w_{t+1})\mid w_t\right]
\le f(w_t)-\eta_t\bigl(1-\tfrac{L\eta_t}{2}\bigr)\norm{\nabla f(w_t)}^2
+\tfrac{L\eta_t^2}{2}\sigma^2 .
\tag{L2}
\]
\end{lemma}
\begin{proof}
\begin{enumerate}
\item Apply Lemma \ref{lem:smooth} with $u=w_{t+1}$, $v=w_t$:
\[
f(w_{t+1})\le f(w_t)+\langle\nabla f(w_t),w_{t+1}-w_t\rangle
+\frac{L}{2}\norm{w_{t+1}-w_t}^2 .
\tag{Step 1}
\]
\item Substitute $w_{t+1}-w_t=-\eta_t g_t$:
\[
f(w_{t+1})\le f(w_t)-\eta_t\langle\nabla f(w_t),g_t\rangle
+\frac{L\eta_t^2}{2}\norm{g_t}^2 .
\tag{Step 2}
\]
\item Take conditional expectation w.r.t.\ $w_t$ and use unbiasedness (A2):
\[
\EE\!\left[\langle\nabla f(w_t),g_t\rangle\mid w_t\right]
= \norm{\nabla f(w_t)}^2 .
\tag{Step 3}
\]
\item Decompose $\norm{g_t}^2$:
\[
\norm{g_t}^2
= \norm{\nabla f(w_t)}^2
+ \norm{g_t-\nabla f(w_t)}^2
+2\langle\nabla f(w_t),g_t-\nabla f(w_t)\rangle .
\tag{Step 4}
\]
\item The cross term vanishes after taking expectation because
$\EE[g_t-\nabla f(w_t)\mid w_t]=0$ (A2). Hence
\[
\EE\!\left[\norm{g_t}^2\mid w_t\right]
= \norm{\nabla f(w_t)}^2+\EE\!\left[\norm{g_t-\nabla f(w_t)}^2\mid w_t\right]
\le \norm{\nabla f(w_t)}^2+\sigma^2 ,
\tag{Step 5}
\]
where the inequality follows from (A3).
\item Plugging (Step 3) and (Step 5) into the expectation of (Step 2) yields
\[
\EE\!\left[f(w_{t+1})\mid w_t\right]
\le f(w_t)-\eta_t\norm{\nabla f(w_t)}^2
+\frac{L\eta_t^2}{2}\bigl(\norm{\nabla f(w_t)}^2+\sigma^2\bigr) .
\tag{Step 6}
\]
\item Rearranging terms gives the claimed inequality (L2). \qedhere
\end{enumerate}
\end{proof}

\newpage

\section*{Main Proof (Step‑by‑Step Derivation)}
We prove Theorem \ref{thm:main} by telescoping the bound from Lemma \ref{lem:descent}.
\begin{enumerate}
\item Take total expectation of Lemma \ref{lem:descent} and drop the conditioning:
\[
\EE\!\left[f(w_{t+1})\right]
\le \EE\!\left[f(w_t)\right]
-\eta_t\Bigl(1-\frac{L\eta_t}{2}\Bigr)\EE\!\left[\norm{\nabla f(w_t)}^2\right]
+\frac{L\eta_t^2}{2}\sigma^2 .
\tag{Step 1}
\]
\item Because of (A4) we have $0<\eta_t\le 1/L$, thus $1-\frac{L\eta_t}{2}\ge \frac12$.
Hence
\[
\EE\!\left[f(w_{t+1})\right]
\le \EE\!\left[f(w_t)\right]
-\frac{\eta_t}{2}\EE\!\left[\norm{\nabla f(w_t)}^2\right]
+\frac{L\eta_t^2}{2}\sigma^2 .
\tag{Step 2}
\]
\item Rearranging (Step 2) gives a lower bound on the gradient term:
\[
\frac{\eta_t}{2}\EE\!\left[\norm{\nabla f(w_t)}^2\right]
\le \EE\!\left[f(w_t)-f(w_{t+1})\right]
+\frac{L\eta_t^2}{2}\sigma^2 .
\tag{Step 3}
\]
\item Sum (Step 3) over $t=0,\dots,T-1$:
\[
\sum_{t=0}^{T-1}\frac{\eta_t}{2}\EE\!\left[\norm{\nabla f(w_t)}^2\right]
\le \EE\!\left[f(w_0)-f(w_T)\right]
+\frac{L\sigma^2}{2}\sum_{t=0}^{T-1}\eta_t^2 .
\tag{Step 4}
\]
\item Since $f(w_T)\ge f^{\star}$,
\[
\EE\!\left[f(w_0)-f(w_T)\right]\le f(w_0)-f^{\star}.
\tag{Step 5}
\]
\item Use the lower bound $\eta_t\ge\eta_{\min}$ and the upper bound $\eta_t\le\eta_{\max}$
to simplify the left‑hand side and the sum of squares:
\[
\sum_{t=0}^{T-1}\frac{\eta_t}{2}\EE\!\left[\norm{\nabla f(w_t)}^2\right]
\ge \frac{\eta_{\min}}{2}\sum_{t=0}^{T-1}\EE\!\left[\norm{\nabla f(w_t)}^2\right],
\tag{Step 6}
\]
\[
\sum_{t=0}^{T-1}\eta_t^2\le T\eta_{\max}^2 .
\tag{Step 7}
\]
\item Plugging (Step 5)–(Step 7) into (Step 4) yields
\[
\frac{\eta_{\min}}{2}\sum_{t=0}^{T-1}\EE\!\left[\norm{\nabla f(w_t)}^2\right]
\le f(w_0)-f^{\star}
+\frac{L\sigma^2}{2}\,T\eta_{\max}^2 .
\tag{Step 8}
\]
\item Divide both sides by $T\eta_{\min}/2$:
\[
\frac{1}{T}\sum_{t=0}^{T-1}\EE\!\left[\norm{\nabla f(w_t)}^2\right]
\le \frac{2\bigl(f(w_0)-f^{\star}\bigr)}{\eta_{\min}T}
+L\eta_{\max}\sigma^2 .
\tag{Step 9}
\]
\item Step (Step 9) is exactly the bound (3) claimed in Theorem \ref{thm:main}.
\end{enumerate}
Thus the theorem follows. \qed

\section*{Rate Derivation (With Constants)}
From (Step 9) we identify two contributions:

\begin{itemize}
\item \textbf{Optimization term:}
\[
C_1:=\frac{2\bigl(f(w_0)-f^{\star}\bigr)}{\eta_{\min}} .
\]
It decays as $C_1/T$.
\item \textbf{Variance term:}
\[
C_2:=L\eta_{\max}\sigma^2 .
\]
It is independent of $T$ and reflects the steady‑state error caused by stochastic noise.
\end{itemize}
Hence
\[
\frac{1}{T}\sum_{t=0}^{T-1}\EE\!\left[\norm{\nabla f(w_t)}^2\right]
\le \frac{C_1}{T}+C_2 .
\]

\section*{Special Cases}
\begin{enumerate}
\item \textbf{Deterministic setting ($\sigma=0$).}
The variance term vanishes and we obtain a pure $O(1/T)$ decay:
\[
\frac{1}{T}\sum_{t=0}^{T-1}\EE\!\left[\norm{\nabla f(w_t)}^2\right]
\le \frac{2\bigl(f(w_0)-f^{\star}\bigr)}{\eta_{\min}T}.
\]
\item \textbf{Constant learning rate.}
If $\eta_{\min}=\eta_{\max}=\eta$, the bound reduces to the classical
non‑convex SGD bound
\[
\frac{2\bigl(f(w_0)-f^{\star}\bigr)}{\eta T}+L\eta\sigma^2 .
\]
\item \textbf{Very long cycles ($M\to\infty$).}
The schedule becomes effectively constant $\eta_t\equiv\eta_{\min}$,
so the bound is governed solely by $\eta_{\min}$; the benefit of a larger
$\eta_{\max}$ disappears.
\end{enumerate}

\end{document}